\appendix

\section{定理 1 的证明}\label{sec:appendix:The1}

本节给出定理 1 中 \emph{Smallest-RLI-First} 策略最优性的形式化证明。

\subsection{问题设置与定义}
我们将 prefill 过程建模为跨层交替出现的通信与计算阶段序列，层编号为 $\ell = 1, \dots, N$。记 $L_{\text{curr}}(t)$ 为时刻 $t$ 正在计算或即将计算的层号。对任意通信流 $f$，记 $L_{\text{target}}(f)$ 为其数据被消费的目标层。相对层索引定义为 $\text{RLI}(f, t) = L_{\text{target}}(f) - L_{\text{curr}}(t)$。我们采用定理中的假设：
\begin{enumerate}
\item[(i)] \textbf{理想计算：} 一旦某层 $\ell$ 的依赖（即所有满足 $L_{\text{target}} = \ell$ 的流）被满足，该层计算便立即开始，并连续运行 $C_\ell$ 时长；
\item[(ii)] \textbf{流体模型：} 通信流可无限细分，因此可无额外开销地进行抢占；
\item[(iii)] \textbf{专用带宽：} 我们聚焦单一链路约束（例如入口或出口瓶颈）。
\end{enumerate}

\subsection{最优性证明}
\textbf{目标。} 最小化 prefill 总工期，等价于最小化最终计算层的完成时间。由于总计算量固定，这又等价于最小化\emph{执行停顿}（即 GPU 因等待数据而空闲的区间）的总时长。

\textbf{区间分类。} 我们依据系统状态，将时间线 $[0, T]$ 划分为两类区间：
\begin{itemize}
\item \textbf{停顿区间（$\mathcal{I}_{\text{stall}}$）：} GPU 因存在某个尚未完成且满足 $\text{RLI}(f, t) = 0$（即目标即为当前层）的流而空闲等待；
\item \textbf{重叠区间（$\mathcal{I}_{\text{overlap}}$）：} GPU 正在执行计算。在这些区间中，链路可用于传输满足 $\text{RLI}(f, t) > 0$ 的流（即未来层所需数据）。
\end{itemize}

\textbf{交换论证。}
假设存在一个最优调度 $\mathcal{S}^*$ 能最小化总工期，但它违反了 Smallest-RLI-First 规则。这意味着在某个时刻 $t \in \mathcal{I}_{\text{stall}}$，调度器把带宽分配给了某条 RLI 更大的流 $y$，而一条 RLI 更小的流 $x$ 仍在等待。

具体而言，由于 $t$ 属于停顿区间，必然存在某条待完成流 $x$ 满足 $\text{RLI}(x) = 0$（即其阻塞当前层）。而违反规则意味着在区间 $[t, t+\delta]$ 内，$\mathcal{S}^*$ 实际服务的是某条流 $y$，且 $\text{RLI}(y) > \text{RLI}(x) = 0$。我们通过如下交换构造一个新调度 $\mathcal{S}'$：
\begin{enumerate}
\item \textbf{交换：} 在 $\mathcal{S}'$ 中，将区间 $[t, t+\delta]$ 分配给流 $x$，而不是流 $y$；
\item \textbf{对停顿的影响：} 由于 $x$ 更早得到服务，当前层 $L_{\text{curr}}$ 的依赖会提前 $\delta$ 个时间单位满足（假设 $x$ 是瓶颈依赖）。这会严格缩短当前 $\mathcal{I}_{\text{stall}}$ 的持续时间，并提前下一段计算阶段（也因此提前下一段 $\mathcal{I}_{\text{overlap}}$）的开始时间；
\item \textbf{$y$ 的可行性：} 流 $y$ 被从时刻 $t$ 挪走。然而，由于 $\text{RLI}(y) > 0$，它直到未来某层才真正需要。$x$ 的完成会触发当前层计算 $C_{L_{\text{curr}}}$ 的开始，从而创造一个新的重叠区间。我们可以安全地把 $y$ 的传输移入这个新产生的重叠区间。由于 $y$ 的截止要求严格晚于 $x$，这种延后不会违反任何依赖关系。
\end{enumerate}

\textbf{结论。} 构造得到的调度 $\mathcal{S}'$ 通过把一部分 $\mathcal{I}_{\text{stall}}$ 转换成 $\mathcal{I}_{\text{overlap}}$，从而减少了累计停顿时长。对所有违反情形重复这一交换论证即可证明：严格优先最小 RLI 的调度（立即服务 $\text{RLI}=0$ 流以最小化停顿，并利用重叠区间服务 $\text{RLI}>0$ 流）能够达到最小总工期。\qed

\section{稳健请求间调度的完整算法} \label{sec:appendix:fullalg}

\begin{algorithm}[t!]
    \caption{请求间调度}
    \small
    \label{alg:ETScheduling_Min}
    \begin{algorithmic}[1]
    \Procedure{InterScheduling}{$\mathcal{B}$, $M$, $B$}
        \FuncInput{$\mathcal{B}, M$: Batch 集合与流量矩阵}
        \FuncInput{$B$: 全局总丢弃预算}
        \FuncOutput{$\sigma, \mathcal{H}$}

        \State $S(\cdot) \gets 0$ \Comment{来自高优先级 batch 的干扰}
        \State $\mathcal{P} \gets \emptyset$; \quad $\mathcal{H} \gets \emptyset$

        \LineComment{步骤 1：按有效截止时间排序}
        \State $\sigma \gets \Call{SortByRED}{\mathcal{B}}$
        \State 从 $M$ 预计算负载向量 $L$

        \For{$k \gets 1$ \textbf{to} $|\sigma|$}
            \State 令 $i \gets \sigma[k]$
            \State $\mathcal{P} \gets \mathcal{P} \cup \mathcal{B}_i$ \Comment{加入候选池}

            \State $\widehat{F}_i \gets \Call{EstFinishTime}{S, L_i}$
            \LineComment{步骤 2：可行性检查}
            \While{$\widehat{F}_i > D^{Lo}_i$ \textbf{and} $|\mathcal{H}| < B$ \textbf{and} $\mathcal{P} \neq \emptyset$}
                \State $u^\star \gets \arg\max_u (S(u) + L_i(u))$ \Comment{瓶颈}
                \LineComment{步骤 3：选择性剪枝}
                \State $r^\star \gets \arg\max_{r \in \mathcal{P}} \text{Load of } r \text{ on } u^\star$

                \State $\mathcal{H} \gets \mathcal{H} \cup \{r^\star\}$; \quad $\mathcal{P} \gets \mathcal{P} \setminus \{r^\star\}$

                \If{$r^\star \in \mathcal{B}_i$} \Comment{丢弃当前 batch 中的请求}
                    \State $L_i(\cdot)\gets L_i(\cdot)-\ell_{r^\star}(\cdot)$
                \Else \Comment{丢弃更高优先级 batch 中的请求}
                    \State $S(\cdot)\gets S(\cdot)-\ell_{r^\star}(\cdot)$
                \EndIf

                \State $\widehat{F}_i \gets \Call{EstFinishTime}{S, L_i}$
            \EndWhile

            \State $S(\cdot) \gets S(\cdot) + L_i(\cdot)$ \Comment{更新}
        \EndFor

        \State \Return $\sigma, \mathcal{H}$
    \EndProcedure
    \end{algorithmic}
\end{algorithm}


算法~\ref{alg:ETScheduling_Min} 展示了整体调度工作流。本节进一步阐释其调度逻辑，并给出优先级分配的完整数学形式以及时延估计与执行机制的更多实现细节。

该节介绍一种面向不精确松弛度与过载场景的稳健调度算法。整个工作流由 batch 到达与离开事件触发，并分为三个阶段。

\parab{步骤 1：通过稳健有效截止时间进行排序。}
为了缓解 “Piggyback Effect”，调度器采用第~\ref{design:interreq} 节定义的稳健有效截止时间（$RED$）指标。优先级计算包含两个阶段。首先，对于包含 $n$ 个请求的 batch $\mathcal{B}$，系统先按截止时间排序，得到 $d_1 \le d_2 \le \dots \le d_n$。随后进行一次线性扫描，找到使相邻截止时间间隔最大的分割点 $k^* = \operatorname*{arg\,max}_{1 \le k < n} (d_{k+1} - d_k)$。这一过程可动态地将离群请求（紧集）与主体请求（松集）区分开来。基于这一划分，调度器结合子 batch 截止时间与离群比例 $f$ 计算 $RED$ 分数。最后，所有活跃 batch 会被组织进一个按照 $RED$ 升序排列的全局优先队列中。这样可以在保留真正紧急工作负载优先级的同时，有效过滤掉由少数离群点造成的短暂“虚高紧迫性”。

\parab{步骤 2：最坏情形可行性检查。} 完成优先级排序后，我们执行准入控制检查，以避免 “Black Hole Effect”。我们在“不同 batch 之间完全不重叠”的最坏假设下，通过严格累积延迟来估计完成时间 $\widehat{T}_i$。计算时延视为确定性，可依据 Transformer 模型的静态计算图和目标硬件上的离线剖析得到~\cite{agrawal2024vidur,distserve,mooncake}。通信时延则通过分析 batch 的流量矩阵，识别网络中的瓶颈端口 $p^*$ 来估计；延迟计算方式为该端口上的累计负载除以其带宽。对于如 Mixture-of-Experts（MoE）这类 token 路由由运行时决定的动态架构，我们依赖历史路由统计。已有工作~\cite{li2025optimizing} 表明，尽管单次流量矩阵噪声较大，但端到端时延预测误差通常仍控制在 20\% 以内。对于我们的调度粒度而言，这样的精度已足够支撑面向目标截止时间（$D^{Lo}_{\min}$）的可行性约束。

\parab{步骤 3：选择性剪枝与软执行。} 如果一个 batch 未通过可行性检查（即 $\widehat{T}_i > D^{Lo}_{\min}$），系统就会触发外科手术式剪枝机制。我们识别出对瓶颈端口 $p^*$ 贡献最大负载的具体请求，并迭代地将其从可行集合中移除，直到预测时延重新满足截止时间。关键的是，这种剪枝采用的是一种\emph{软执行}策略。调度器不会立刻丢弃被剪枝请求，而是将它们降级到一个“双队列系统”中的 “Scavenger” 优先级类别。主队列（Main Queue）采用严格优先服务，而 Scavenger 队列则只在检测到网络空闲、或实际运行时延低于悲观估计时机会性地被服务。该混合策略既能在最坏预测条件下维持系统稳定，又能在拥塞并未真正发生时回收资源、尽可能提高吞吐。

\parab{总结。} 需要强调的是，这一工作流内部具有清晰的设计层次。RED 指标是首要的决策依据，它控制全局发放顺序，在正常到中等负载下最大化服务级目标（SLO）满足率。可行性检查与选择性剪枝则仅作为互补性的保险机制，在病态过载场景下启动，以防止资源耗尽。这样的职责分离确保了调度器在绝大多数情况下保持稳定、可预测——主要依赖稳健排序来做决策——同时在极端争用不可避免时，仍具备平滑降级的能力。
